\scanpage{106}
\mspage{342}

(VII) Calculations in $\mathrm{Sl}(2,\widehat{\mathbb Z})'=\mathrm{Sl}(2,\widehat{\mathbb Z})/\pm1$
(\unsure{cocycles} under $\rho$).

\margin{Calculations ``for conscience's sake'', \unsure{which will not be of use} perhaps any more!}

This group is a (very small!) quotient of $\widehat{\mathrm{Sl}(2,\mathbb Z)'}$, which has the advantage of \unsure{lending} itself to convenient explicit \inserted{algebraic} calculations, in terms of matrices. Let us denote, by abuse of language, by the same letters $g,g_0,h_0,\xi,\lambda_0,\lambda$ elements of \deleted{\illegible{1}} $\widehat{\Pi}_{0,3}$, \inserted{(or of $\widehat{\mathfrak S}_{0,3}^{+}$ more generally)} and their images in \inserted{$\Pi(2)'\subset$} $\mathrm{Sl}(2,\widehat{\mathbb Z})'$ \deleted{\illegible{1}} (or in $\mathrm{Gl}(2,\widehat{\mathbb Z})'\simeq\mathrm{Gl}(2,\widehat{\mathbb Z})/\{\pm1\}$). \deleted{We} \deleted{\unsure{shall set}} in general, let us consider a matrix
\[
(76)\qquad g=\begin{pmatrix}a&b\\c&d\end{pmatrix}\in\mathrm{Gl}(2,\widehat{\mathbb Z})
\]
we then have
\[
(77)\quad
\left\{\begin{aligned}
\deleted{\unsure{\rho}}\rho(g)&=\rho g\rho^{-1}=\begin{pmatrix}c+d&-c\\c+d-a-b&a-c\end{pmatrix},
&\inserted{\det(g)}\rho(g^{-1})&=\begin{pmatrix}a-c&c\\a+b-c-d&c+d\end{pmatrix}\\
\rho^2(g)&=\rho^{-1}g\rho=\deleted{\unsure{(}}\begin{pmatrix}d-b&a+b-c-d\\-b&a+b\end{pmatrix},
&\inserted{\det(g)}\rho^2(g^{-1})&=\begin{pmatrix}a+b&c+d-a-b\\\deleted{\unsure{c}}b&d-b\end{pmatrix}
\end{aligned}\right.
\]
\[
(78)\quad
\left\{\begin{aligned}
\relax[g,\rho]&=g\rho g^{-1}\rho^{-1}=\begin{pmatrix}A&B\\C&D\end{pmatrix}.\ \deleted{\unsure{\frac{1}{\det(g)}}}\\
\delta A&=(a^2+b^2+ab)-(ac+bd+bc)&\delta B&=ac+bd+bc\\
\delta C&=-(c^2+d^2+cd)+(ac+bd+ad)&\delta D&=c^2+d^2+cd
\end{aligned}\right.
\]
\margin{N.B. $C+D-B=\unsure{1}\deleted{\unsure{\det g=\delta}}$\quad $\delta=\det g=ad-bc$ \deleted{\illegible{1}}}

We shall write the equation
\[
\deleted{(\unsure{79})}\qquad g\rho(g)\rho^2(g)=1
\]
in the equivalent form
\[
(79)\qquad g\rho(g)=\rho^{-1}(g^{-1})\bigl(=\rho^2(g^{-1})\bigr)
\qquad\text{\uline{in $\mathrm{Gl}(2,\widehat{\mathbb Z})'$}}
\]
where
\[
(80)\qquad g\rho(g)=\begin{pmatrix}P&Q\\R&S\end{pmatrix}
\]

\scanpage{107}

\[
(81)\quad
\left\{\begin{aligned}
P&=(a+b)(c+d-b)&Q&=-a(c-b)-bc\overset{\inserted{\deleted{\text{if }\det g=+1}}}{=}\deleted{\unsure{1}}\delta-a(c+d-b)\\
R&=(c+d)^2-d(a+b)&S&=ad-c(c+d)=\underset{\inserted{\text{if }\det g=1}}{\delta}-c(c+d-b)\\
&\underset{\inserted{\text{if }\det g=+1}}{=}(c+d)(c+d-b)-\delta
\end{aligned}\right.
\]
where we have set $\delta=ad-bc=\det g$.

Comparing (81) and the expression (77) for $\rho^2(g^{-1})$, equation (79) -- which is also written
\[
g\rho(g)=\pm\rho^2(g^{-1})
\]
gives the following solutions.

\margin{Up to replacing $g$ by $\unsure{-g}$, \deleted{\illegible{5}} \unsure{we are reduced} to \unsure{determining} the solutions of the equation}

\[
(82)\qquad \unsure{g|\rho}\rho(g)\rho^2(g)=\unsure{+|\pm}1\qquad\text{i.e.}\quad \unsure{g|\rho}\rho(g)=\unsure{+|\pm}\rho^{-1}(g^{-1})\qquad\text{\uline{in $\mathrm{Gl}(2,\widehat{\mathbb Z})$}}
\]
\deleted{$\unsure{(82)}$ \illegible{2} $\unsure{g\rho(g)=+\rho^{-1}(g^{-1})}$.}
We distinguish the solutions according to the factors of the product
\[
\mathrm{Sl}(2,\widehat{\mathbb Z})'\simeq\prod_{\ell\ \mathrm{premier}}\mathrm{Sl}(2,\mathbb Z_\ell)'.
\]
\margin{\uline{N.B.} \unsure{We} \unsure{shall use|use} this study of equation (82) only in the case $\det g=1$, i.e. $\delta=1$, \unsure{unless I am mistaken}}

The relation
\[
\frac1\delta(a+b)=P=(a+b)(c+d-b)\qquad\text{in }\mathbb Z_\ell
\]
gives the two solutions $a+b=0$ or $c+d-b=1$ (since $\mathbb Z_\ell$ is an integral domain).

\deleted{1°) $(a+b)\ne0$ (in $\mathbb Z_\ell$) whence $c+d-b=1$
\[
(83)\qquad\boxed{c+d-b=\frac1\delta}.
\]
Given (83), we verify \unsure{at once} that the formulas \inserted{$g\rho(g)=\rho^{-1}(g^{-1})$} -- which remain to be checked
\[
\begin{aligned}
P&=(a+b)/\delta&Q&=(1-a)/\delta\\
R&=b/\delta&S&=(1-c)/\delta
\end{aligned}
\]
\inserted{\deleted{\illegible{1}} $\unsure{R=}$
\[
\begin{aligned}
Q&=\frac{c+d-\unsure{a}-b}{\delta}\left(=\unsure{\frac1{\delta^2}|\frac1\delta}-\frac a\delta\right),\\
R&=\frac b\delta,\quad S=\frac{d-b}{\delta}\left(=\frac1{\delta^2}-\frac c\delta\right)
\end{aligned}
\]}
(84) are equivalent to \illegible{5} (82) \illegible{1}
\[
(84)\qquad\boxed{\delta^3=1}.
\]
}

\margin{The necessity of this formula (84) was obvious a priori.}

\deleted{2°) $a+b=0$ i.e. $a=-b$, i.e. $g=\begin{pmatrix}-b&b\\c&d\end{pmatrix}$, and
\[
-b(c+d)=1
\]
\[
(84)\quad\left\{\begin{aligned}
\det g&=-b(c+d)=\delta\quad(\text{thus }\delta+b(c+d)=0;\unsure{)}\quad d=-\frac\delta b-c\\
P&=0,&Q&=\delta+b(c+d-b)=-b^2\\
R&=(c+d)^2=\frac{\delta^2}{b^2},&S&=\unsure{ad}\ \illegible{1}\ {-c(c+d)+ad=1-c(\unsure{c+d-b})}\\
&&&\deleted{\illegible{1}c}=\delta+c\left(\frac\delta b+b\right)
\end{aligned}\right.
\]
}

\scanpage{108}
\mspage{343}

1°) $a+b\ne0$ (in $\mathbb Z_\ell$) whence
\[
(83)\qquad\boxed{c+d-b=\frac1\delta}.
\]
Given (83), we verify that the formulas (82), -- \unsure{where} it remains to check
\[
\begin{aligned}
Q&=\frac{c+d-(a+b)}\delta\ \unsure{\leftarrow|=}\left(=\frac1{\delta^2}-\frac a\delta\right),&R&=\frac b\delta\\
S&=\frac{d-b}\delta.\ \left(=\frac1{\delta^2}-\frac c\delta\right)
\end{aligned}
\]
are equivalent to:
\[
(84)\qquad\boxed{\delta^3=1}
\]
(whose necessity was obvious a priori).

\margin{\uline{N.B.} The conditions $c+d-b=\frac1\delta$, $\delta^3=1$ are sufficient \illegible{1} for equations (82), \unsure{!} (\unsure{independently} of the question whether $a+b=0$ or $a+b\ne0$) -- they are necessary \unsure{if we assume that we} \unsure{are not} in case 2° below}

2°) $a+b=0$ (in $\mathbb Z_\ell$) whence
\[
(85)\qquad a=-b\qquad\text{i.e.}\quad g=\begin{pmatrix}-b&b\\c&d\end{pmatrix}
\]
\[
(86)\quad\left\{\begin{aligned}
\det g&=-b(c+d)=\delta&\text{whence}\\
\deleted{\unsure{b+}\delta}\ \delta+b(c+d)&=0&d&=-\frac\delta b-c
\end{aligned}\right.
\]
i.e.
\[
g=\begin{pmatrix}-b&b\\c&-\frac\delta b-c\end{pmatrix},
\]
and
\[
\left\{\begin{aligned}
P&=0,&Q&=\delta+b(c+d-b)=-b^2\\
R&=(c+d)^2=\frac{\delta^2}{b^2},&S&=\delta+c\left(\frac\delta b+b\right);
\end{aligned}\right.
\]
\deleted{Equation (82) is then written, since \unsure{here}}
\[
\rho^2(g^{-1})=\begin{pmatrix}0&-1/b\\b/\delta&-\frac1\delta\left(\frac\delta b+b+c\right)\end{pmatrix},
\]
\deleted{and and}

and equation (82) is then written
\[
(87)\quad\left\{\begin{aligned}
\text{a)}\quad&b^3=\delta^3=1.\\
\text{b)}\quad&c\bigl((b\delta)^2+(b\delta)+1\bigr)+(\delta^2+b^2\delta+b)=0.
\end{aligned}\right.
\]

\scanpage{109}

We again distinguish two cases:

a) $b\delta=1$ i.e. $b=\frac1\delta$, then equation \inserted{87 b)} becomes
\[
3c+3\frac1\delta=0\qquad\text{i.e.}
\]
\deleted{$3c=$}
\[
(88)\qquad c=-\frac1\delta=\inserted{-\delta^2}\ \deleted{\unsure{3c+(2\delta^2+1)=0}}\qquad\deleted{\unsure{3c+3\delta^2=0}\quad\text{i.e.}}
\]
\deleted{Now $\delta\in\widehat{\mathbb Z}^{*}\Rightarrow\delta^2\equiv1\ (3)$ thus $2\delta^2+1\equiv0\ (3)$, whence the existence of a solution in $\widehat{\mathbb Z}$ of (88)
\[
(89)\qquad c=-\frac13(2\delta^2+1)
\]
}
corresponds to
\[
(90)\qquad g=\deleted{\begin{pmatrix}
-\frac1\delta=-\delta^2&\frac1\delta=\delta^2\\
\unsure{-\frac13(2\delta^2+1)}\ \inserted{-\delta^2}&\frac13(\unsure{4|1}-\delta^2)
\end{pmatrix}}\quad
\begin{pmatrix}
-\delta'\deleted{\unsure{{}^{2}}}&\delta'\deleted{\unsure{{}^{2}}}\\
-\delta'\deleted{\unsure{{}^{2}}}&0
\end{pmatrix}
\]
\inserted{where $\delta'=1/\delta=\delta^2$, and} where, I recall, $\delta$ is subject to the sole condition
\[
\delta^3=1.
\]

b) $b\delta\ne1$ thus (since $(b\delta)^3=1$) $(b\delta)^2+b\delta+1=0$)
and equation (87 b) becomes
\[
(\alpha)\qquad\delta^2+b^2\delta+b=0.\qquad\deleted{\text{\unsure{i.e. a}}}
\]
We shall have
\[
\left\{\begin{aligned}
b&=\deleted{\unsure{\delta}}u\frac1\delta=u\delta^2\quad\text{with}\quad u^3=1,\ u\ne1\quad\deleted{\text{thus}}\quad\text{i.e.}\\
u^2+u+1&=0
\end{aligned}\right.
\]
and equation $\alpha)$, \deleted{is written} reduces to
\[
\unsure{u^2|\delta^2}+u\deleted{\unsure{{}^{2}}}+1=0.
\]
\deleted{\illegible{1}}

\deleted{which requires $\delta\ne1$ (otherwise $u^2=0$, \deleted{\illegible{1}} incompatible with $u^3=1$) thus $\delta^2+\delta+1=0$ -- thus $(\alpha)$ is written
\[
-(\delta+1)-(u+1)+1=0\qquad\text{i.e.}
\]
\[
\delta+u+1=0
\]
which \deleted{is \unsure{obviously} \illegible{3}} \inserted{\unsure{means}} that $u,\delta$ are \uline{distinct} solutions of the cyclotomic equation $x^2+x+1=0$. Now $u\ne\delta$ \uline{means} \unsure{also that} $\unsure{\frac b\delta\ne\delta}$ i.e. $b\ne\delta^2$ \deleted{which requires} i.e. $b=\delta$.
\[
(91)\qquad g=\begin{pmatrix}-\delta&\delta\\c&-1-c\end{pmatrix}
\]
}

\scanpage{110}
\mspage{344}

\uline{Proposition} Let $g\in\mathrm{Gl}(2,\widehat{\mathbb Z}_\ell)$, $g=\begin{pmatrix}a&b\\c&d\end{pmatrix}$, $a,b,c,d\in\unsure{\widehat{\mathbb Z}_\ell|\mathbb Z_\ell}$,
$ad-bc=\delta\in\unsure{\widehat{\mathbb Z}_\ell^{*}|\mathbb Z_\ell^{*}}$. For us to have
\[
g\rho(g)\rho^2(g)=1\qquad(\unsure{82|81})
\]
it is necessary and sufficient that we have $\delta^3=1$, and that we be in one of the following two cases, which are mutually exclusive:
\[
\begin{aligned}
1\text{°)}\quad&c+d-b=1\\
2\text{°)}\quad&g=\deleted{\left(\begin{matrix}-\delta'&\delta'\end{matrix}\right.}\ -\delta'\begin{pmatrix}-1&1\\1&0\end{pmatrix}=-\delta'\rho^{-1}\qquad\left(\text{where }\delta'=\frac1\delta=\delta^2\right)
\end{aligned}
\]

\deleted{Moreover, in case 1°), we have $a+b=0$ \unsure{iff|sss} \unsure{when of}}

\uline{Corollary} The $g$ for which we have \deleted{at the same \uuline{time}} \inserted{1°) and} $a+b\ne0$ are those of the form
\[
\deleted{g=\left(\begin{matrix}-\unsure{u}\delta'&u\delta'\\\unsure{c|0}&\end{matrix}\right.}
\]
\[
\left\{\begin{aligned}
g&=\begin{pmatrix}a&b\\c&b-c+\frac1\delta\end{pmatrix}\qquad\deleted{\unsure{\text{for}}\ \unsure{c}=\frac1{(a+b)}\unsure{(a(b+\delta^2))|(ab+a\delta')}}\\
&\text{with}\quad c=\frac1{a+b}\bigl(a(b+\delta')-\delta\bigr),
\end{aligned}\right.
\]
those for which we have 1°) and $a+b=0$ are \deleted{those} the matrices of the form
\[
\left\{\begin{aligned}
g&=\begin{pmatrix}-u\delta'&u\delta'\\c&-\delta'u'-c\end{pmatrix}\\
&\text{where }\delta'=1/\delta=\delta^2,\ u'=\frac1u=u^2,\quad u\text{ being a solution}\\
&\text{of the cyclotomic equation }u^2+u+1=0\quad\text{i.e.}\\
&\text{\unsure{root|roots} \uline{\unsure{primitive|primitive}} }3^{\mathrm e}\text{ of unity}
\end{aligned}\right.
\]